% Minimalist adaptation, 2026. See ATTRIBUTION.md for provenance and changes.
% Independent original example by the schematic-designer maintainers.
% Pattern reference (no upstream code copied):
% https://github.com/f0nzie/tikz_favorites/blob/9506fd5643abb406197bed176ab9183484006dd1/src/plot-linear-regression%2Bgeometry.tex
% License: MIT, as part of the biophysics-simulation-workflow-skills bundle.
\documentclass[tikz,border=8bp]{standalone}
\usepackage{minimalist-profile}
\usetikzlibrary{arrows.meta,calc}

\begin{document}
\begin{tikzpicture}[ms base,x=1.15cm,y=1.15cm]
  \pgfmathsetmacro{\xP}{1.35}
  \pgfmathsetmacro{\xQ}{4.85}
  \pgfmathsetmacro{\yP}{0.16*(\xP-2.65)^2+0.90}
  \pgfmathsetmacro{\yQ}{0.16*(\xQ-2.65)^2+0.90}

  \coordinate (O) at (0,0);
  \coordinate (P) at (\xP,\yP);
  \coordinate (Q) at (\xQ,\yQ);
  \coordinate (Px) at (\xP,0);
  \coordinate (Qx) at (\xQ,0);
  \coordinate (R) at (\xQ,\yP);

  \draw[-{Latex[length=\msArrowLength,width=\msArrowWidth]},draw=msInk,line width=\msLineStructure]
    (-0.2,0) -- (6.35,0) node[below] {$x$};
  \draw[-{Latex[length=\msArrowLength,width=\msArrowWidth]},draw=msInk,line width=\msLineStructure]
    (0,-0.2) -- (0,3.25) node[left] {$f(x)$};

  \draw[color=msColor1,line width=\msLineEmphasis]
    plot[domain=0.15:6.05,samples=100,smooth]
      (\x,{0.16*(\x-2.65)^2+0.90});
% Function label matches the analytic curve; secant has a separate semantic color.
  \node[anchor=west,ms body,text=msColor1] at (6.05,2.9) {$f$};

  \draw[color=msColor5,line width=\msLineStructure]
    ($(P)!-0.22!(Q)$) -- ($(P)!1.22!(Q)$)
    node[pos=0.75,above,sloped,ms body,text=msColor5] {secant};

  \draw[ms context,draw=msMuted,ms dashed] (P) -- (Px);
  \draw[ms context,draw=msMuted,ms dashed] (Q) -- (Qx);
  \draw[ms context,draw=msMuted,ms dashed] (P) -- (R) -- (Q);

  \fill[msPaper,draw=msInk,line width=\msLineStructure] (P) circle ({0.5*\msMarkerPoint})
    ;
  \node[ms body,text=msInk,anchor=south] at (1.35,1.5) {$P=(x_0,f(x_0))$};
  \fill[msPaper,draw=msInk,line width=\msLineStructure] (Q) circle ({0.5*\msMarkerPoint})
    ;
  \node[ms body,text=msInk,anchor=east] at (4.75,0.55) {$Q=(x_0+\Delta x,f(x_0+\Delta x))$};

  \draw[{Latex[length=\msArrowLength,width=\msArrowWidth]}-{Latex[length=\msArrowLength,width=\msArrowWidth]},color=msInk]
    ([yshift=-\msSpaceGroup]Px) -- node[below,ms body,text=msInk] {$\Delta x$} ([yshift=-\msSpaceGroup]Qx);
  \draw[{Latex[length=\msArrowLength,width=\msArrowWidth]}-{Latex[length=\msArrowLength,width=\msArrowWidth]},color=msInk]
    ([xshift=\msSpaceGroup]R) -- node[right,ms body,text=msInk] {$\Delta f$} ([xshift=\msSpaceGroup]Q);

  \node[ms equation,inner sep=\msSpaceLabel,anchor=north west] at (0.35,3.02)
    {$m_{\mathrm{sec}}=\dfrac{\Delta f}{\Delta x}
      =\dfrac{f(x_0+\Delta x)-f(x_0)}{\Delta x}$};
\end{tikzpicture}
\end{document}
